\documentclass[12pt]{article}
\usepackage[utf8]{inputenc}
\usepackage[T1]{fontenc}
\usepackage{hyperref}
\title{Shihhi Arabic Language Resources for K-12 Curriculum Development}
\author{Compiled by Claude Sonnet 5.5 (claude-sonnet-5-5) for Andr\'as Kornai}
\date{}
\begin{document}
\maketitle

\section{Introduction}
Shihhi Arabic (ISO 639-3: ssh; Glottocode shih1239) is an Arabic variety spoken by the Shihuh in the Musandam Governorate of Oman and in Ras Al Khaimah, United Arab Emirates, with some speakers in Iran. Wikipedia gives about 38,000 speakers (2020) and Global Recordings Network about 39,000. It should not be confused with Kumzari, an Iranian language spoken by the Kumzarah subgroup of the Shihuh. Search results showed very little documentation or teaching material specific to Shihhi.

\section{Primary Educational Resources}
No Shihhi primers, textbooks or school programmes were found. Children in the region are presumably taught in Modern Standard Arabic, but this was not confirmed in the search results. A curriculum would need to treat Shihhi as a home variety alongside Standard Arabic.

\section{Linguistic and Reference Materials}
\begin{itemize}
\item Wikipedia article on Shihhi Arabic. It mentions a possible inland and coastal dialect split and cites a 2011 Leiden University thesis by Roy S. Bernabela, a sketch of the phonology and morphology of the al-Jadida dialect in Oman. I did not verify the thesis itself: \url{https://en.wikipedia.org/wiki/Shihhi_Arabic}
\item Wikipedia article on the Shihuh people: \url{https://en.wikipedia.org/wiki/Shihuh}
\end{itemize}

\section{Religious and Community Resources}
Global Recordings Network has a language page for Shihhi, which reports that no audio recordings are currently available: \url{https://globalrecordings.net/en/language/ssh}

\section{Digital and Audio Resources}
\begin{itemize}
\item Aletihad news report on a Shihuh researcher working on the oral heritage of the Shihuh in Ras Al Khaimah, interviewing elders. The page date looks unreliable, and the URL as returned is odd: \url{https://en.aletihad.ae/news/culture/4688139/javasxript:}
\item Detect Language entry for Shihhi Arabic (a reference listing only): \url{https://detectlanguage.com/languages/ssh}
\end{itemize}

\section{Curriculum Development Context}
Documentation is thin: one thesis on the Omani side, no audio archive and some community oral-history work in Ras Al Khaimah. A curriculum would first need basic documentation, such as word lists, recorded texts and a decision on how to write the variety. That work should involve Shihuh communities in Ras Al Khaimah and Musandam. Dictionaries and corpora were not searched for beyond the queries above, so more material may exist.

\end{document}